\chapter{Graphs, Networks, and Markov Chains}

So far in this book the vectors and matrices have mostly described \emph{geometry}:
arrows in the plane, stretches of the unit square, rotations of space. In this chapter
we point the very same machinery at something that does not, at first, look geometric at
all: a \emph{network}. A web of friendships. A power grid. The links between web pages.
The way a random walker hops from city to city. A spreadsheet of who emails whom.

The remarkable thing is that the structure of a network is captured \emph{exactly} by a
matrix, and the deep questions you might ask about the network (Is it connected? Where are
its natural communities? Which page is the most important? Where does a random walker end
up in the long run?) turn into questions about the \emph{eigenvalues and eigenvectors} of
that matrix. The graph is the picture; the matrix is the algebra; the spectrum is the
hidden geometry. Linear algebra is the bridge.

By the end of this chapter you will be able to look at a graph and write down three
matrices for it (incidence, adjacency, Laplacian); you will know what each row, column,
and eigenvalue is telling you; and you will understand PageRank -- the algorithm that
launched a trillion-dollar company -- as nothing more than finding an eigenvector.

\section{From pictures to matrices}
\label{ch35:pictures-to-matrices}

A \textbf{graph} is the simplest object in all of mathematics that still has interesting
structure. It is a set of \emph{nodes} (also called \emph{vertices}) and a set of
\emph{edges}, where each edge joins two nodes. That is the whole definition. Nodes are
dots; edges are the lines between them. Draw five dots, connect some of them with lines,
and you have a graph.

\begin{intuitionbox}
A graph throws away everything except \emph{who is connected to whom}. Where you place
the dots on the page does not matter -- only the pattern of connections does. Two drawings
that look completely different can be the ``same'' graph if the connection pattern matches.
The job of linear algebra is to turn that bare connection pattern into numbers we can
compute with, and then to recover meaning from those numbers.
\end{intuitionbox}

Let us fix a small running example we will use throughout. It has four nodes labeled
$1,2,3,4$ and five edges. Figure~\ref{ch35:fig-graph4} shows it.

\begin{figure}[H]
\centering
\begin{tikzpicture}[scale=1.6,
  node/.style={circle,draw=sourceblue,very thick,fill=blue!8,minimum size=8mm,inner sep=0pt,font=\small},
  every edge/.style={draw=spacegray,very thick}]
  \node[node] (n1) at (0,1.3) {$1$};
  \node[node] (n2) at (1.5,1.3) {$2$};
  \node[node] (n3) at (1.5,0) {$3$};
  \node[node] (n4) at (0,0) {$4$};
  \draw (n1) edge (n2);
  \draw (n1) edge (n4);
  \draw (n2) edge (n3);
  \draw (n2) edge (n4);
  \draw (n3) edge (n4);
\end{tikzpicture}
\caption{Our running graph $G$: four nodes, five edges. Node $2$ touches three edges
(it has \emph{degree} $3$); node $1$ touches two. There is no edge between $1$ and $3$.}
\label{ch35:fig-graph4}
\end{figure}

The \emph{degree} of a node is the number of edges touching it. In our graph the degrees
are: node $1$ has degree $2$, node $2$ has degree $3$, node $3$ has degree $2$, node $4$
has degree $3$. A useful sanity check you should burn into memory: the degrees sum to
twice the number of edges, because every edge contributes $1$ to the degree of each of its
two endpoints. Here $2+3+2+3 = 10 = 2 \times 5$. Good.

\subsection{The adjacency matrix}
\label{ch35:adjacency}

The most direct way to record a graph as numbers is the \textbf{adjacency matrix} $\mat{A}$.
For a graph with $n$ nodes it is the $n \times n$ matrix whose $(i,j)$ entry is

\[
A_{ij} =
\begin{cases}
1 & \text{if there is an edge between node $i$ and node $j$,}\\
0 & \text{otherwise.}
\end{cases}
\]

For our graph, reading off Figure~\ref{ch35:fig-graph4} row by row:

\[
\mat{A} =
\begin{bmatrix}
0 & 1 & 0 & 1\\
1 & 0 & 1 & 1\\
0 & 1 & 0 & 1\\
1 & 1 & 1 & 0
\end{bmatrix}.
\]

 Look at row $2$: it reads $1,0,1,1$, meaning node $2$ connects to nodes $1$, $3$, and $4$
but not to itself. The diagonal is all zeros (no node has an edge to itself), and the
matrix is \emph{symmetric}, $\mat{A} = \mat{A}\T$, because ``$i$ connects to $j$'' is the
same statement as ``$j$ connects to $i$.'' That symmetry is the algebraic shadow of the edges being \emph{undirected}, and it is what
will later guarantee real eigenvalues.

\begin{focusbox}
Three facts about an adjacency matrix of an undirected graph, worth holding onto:
\begin{itemize}
\item It is symmetric: $\mat{A} = \mat{A}\T$.
\item Its diagonal is zero (for a simple graph with no self-loops).
\item The sum of row $i$ equals the degree of node $i$. Here the row sums are
$2,3,2,3$ -- exactly the degree sequence.
\end{itemize}
\end{focusbox}

Why store a graph this way? Because once it is a matrix, \emph{matrix operations have
graph meaning}. The first example of this is already a striking one.

\begin{intuitionbox}
\textbf{Powers of $\mat{A}$ count walks.} The entry $(\mat{A}^k)_{ij}$ is exactly the
number of walks of length $k$ from node $i$ to node $j$ (a walk is a sequence of edges,
repeats allowed). Multiplying $\mat{A}$ by itself is the same as ``take one more step
along an edge,'' and matrix multiplication adds up all the ways to do that.
\end{intuitionbox}

Let us verify the first nontrivial case. The $(i,j)$ entry of $\mat{A}^2$ is
$\sum_k A_{ik}A_{kj}$. The product $A_{ik}A_{kj}$ is $1$ exactly when there is an edge
$i\!-\!k$ \emph{and} an edge $k\!-\!j$, i.e. when $k$ is a common neighbor giving a length-2
walk $i \to k \to j$. Summing over $k$ counts all such intermediate nodes. So
$(\mat{A}^2)_{ij}$ counts length-$2$ walks. Let us compute it for our graph:

\[
\mat{A}^2 =
\begin{bmatrix}
0&1&0&1\\1&0&1&1\\0&1&0&1\\1&1&1&0
\end{bmatrix}
\begin{bmatrix}
0&1&0&1\\1&0&1&1\\0&1&0&1\\1&1&1&0
\end{bmatrix}
=
\begin{bmatrix}
2&1&2&1\\1&3&1&2\\2&1&2&1\\1&2&1&3
\end{bmatrix}.
\]

Read the diagonal: $(\mat{A}^2)_{22} = 3$. That says there are three length-$2$ walks from
node $2$ back to node $2$ -- namely out to each of its $3$ neighbors and back. In general
$(\mat{A}^2)_{ii}$ is just the degree of node $i$. The off-diagonal entry
$(\mat{A}^2)_{13} = 2$ says there are two length-$2$ paths from $1$ to $3$: through node
$2$, and through node $4$. Trace your finger over Figure~\ref{ch35:fig-graph4} to confirm
it.

\begin{labbox}
Compute $\mat{A}^3$ for the graph above and read off $(\mat{A}^3)_{12}$. It should equal
the number of length-$3$ walks from node $1$ to node $2$. Then find them by hand on the
picture (start at $1$, take three edges, end at $2$). Do the two answers agree? This little
exercise is the whole reason adjacency matrices earn their keep.
\end{labbox}

\subsection{Directed graphs and the adjacency matrix that is not symmetric}

If edges have a \emph{direction} -- web links, one-way streets, ``follows'' on a social
network, the flow of money -- we set $A_{ij}=1$ when there is an arrow \emph{from} $i$
\emph{to} $j$. Now $\mat{A}$ need not be symmetric: page $i$ can link to page $j$ without
$j$ linking back. We will need exactly this asymmetry for PageRank at the end of the
chapter. For now, just notice that the language ``$\mat{A}=\mat{A}\T$'' is precisely the
dividing line between undirected and directed.

\subsection{The incidence matrix: edges meet nodes}
\label{ch35:incidence}

The adjacency matrix is node-versus-node. There is a second, equally important matrix that
is \emph{edge}-versus-node, and it is the one electrical engineers and Gil Strang reach for
first. To build it, orient each edge -- pick an arbitrary ``from'' and ``to''
for every edge, just to fix signs. (For an undirected graph the orientation is a
bookkeeping choice; the final physics will not depend on it.)

Number the five edges of our graph and orient them:
$e_1: 1\!\to\!2$, $\;e_2: 1\!\to\!4$, $\;e_3: 2\!\to\!3$, $\;e_4: 2\!\to\!4$,
$\;e_5: 3\!\to\!4$. The \textbf{incidence matrix} $\mat{B}$ has one row per edge and one
column per node. In row $e$, we put $-1$ in the ``from'' column, $+1$ in the ``to''
column, and $0$ everywhere else:

\[
\mat{B} =
\bordermatrix{%
      & 1 & 2 & 3 & 4 \cr
e_1\; & -1 & +1 & 0 & 0 \cr
e_2\; & -1 & 0 & 0 & +1 \cr
e_3\; & 0 & -1 & +1 & 0 \cr
e_4\; & 0 & -1 & 0 & +1 \cr
e_5\; & 0 & 0 & -1 & +1 \cr}.
\]

Each row has exactly one $-1$ and one $+1$: an edge touches exactly two nodes. This little
matrix is a difference operator in disguise.

\begin{intuitionbox}
Think of a column vector $\vec{x} = (x_1,x_2,x_3,x_4)\T$ as assigning a number to each
node -- a \emph{potential}, a temperature, an elevation. Then $\mat{B}\vec{x}$ assigns to
each \emph{edge} the difference of the potentials at its two ends. Row $e_1$ of
$\mat{B}\vec{x}$ is $x_2 - x_1$. So $\mat{B}$ is the discrete ``gradient'': it turns
node-values into edge-differences. This is the graph version of $\frac{d}{dx}$.
\end{intuitionbox}

Let us make that concrete. Take $\vec{x} = \colvec{5\\8\\6\\2}$ (an elevation at each node).
Then

\[
\mat{B}\vec{x} =
\begin{bmatrix}
-1&1&0&0\\-1&0&0&1\\0&-1&1&0\\0&-1&0&1\\0&0&-1&1
\end{bmatrix}
\colvec{5\\8\\6\\2}
=
\colvec{8-5\\2-5\\6-8\\2-8\\2-6}
=
\colvec{3\\-3\\-2\\-6\\-4}.
\]

Each entry is the rise (or fall) along that edge. If $\vec{x}$ were a voltage, these would
be the voltage \emph{drops} across each wire -- which, divided by resistance, give the
currents. The incidence matrix is the direct algebraic encoding of Kirchhoff's laws.

\begin{strangbox}
If you have read Strang, you will recognize this as his beloved framework
$\mat{A}\T \mat{C} \mat{A}$ from the chapter on graphs and networks (he calls the incidence
matrix $\mat{A}$). The four subspaces of the incidence matrix tell the whole story: its
\emph{nullspace} is the constant vectors (adding the same elevation everywhere changes no
difference), and its \emph{left nullspace} is the loops in the graph (Kirchhoff's voltage
law -- drops around any loop sum to zero). If you are reading this \emph{before} Strang,
just hold the picture ``$\mat{B}$ takes differences across edges''; when you meet his
$\mat{A}\T \mat{A}$, you will already know it is the Laplacian we are about to build.
\end{strangbox}

What lives in the nullspace of $\mat{B}$? We need $\mat{B}\vec{x} = \vec{0}$, i.e. every
edge has zero difference, i.e. $x_i = x_j$ for every edge $i\!-\!j$. If the graph is
connected, that forces \emph{all} the $x_i$ to be equal: the nullspace is exactly the
multiples of $\vec{1} = (1,1,1,1)\T$, a one-dimensional space. This is the discrete version
of ``a function with zero gradient is constant.'' Keep this in your pocket -- it returns in
two pages with deep meaning.

\section{The graph Laplacian}
\label{ch35:laplacian}

We now build the single most important matrix attached to a graph. Combine the degrees and
the adjacency into the \textbf{graph Laplacian}:

\begin{definition}[Graph Laplacian]
\label{ch35:def-laplacian}
Let $\mat{D} = \diag(d_1,\dots,d_n)$ be the diagonal \emph{degree matrix}, where $d_i$ is
the degree of node $i$, and let $\mat{A}$ be the adjacency matrix. The graph Laplacian is
\[
\mat{L} = \mat{D} - \mat{A}.
\]
Equivalently, $L_{ii} = d_i$, and for $i \ne j$, $L_{ij} = -1$ if $i,j$ are adjacent and
$0$ otherwise.
\end{definition}

For our running graph, $\mat{D} = \diag(2,3,2,3)$, so

\[
\mat{L} = \mat{D} - \mat{A} =
\begin{bmatrix}2&0&0&0\\0&3&0&0\\0&0&2&0\\0&0&0&3\end{bmatrix}
-
\begin{bmatrix}0&1&0&1\\1&0&1&1\\0&1&0&1\\1&1&1&0\end{bmatrix}
=
\begin{bmatrix}
2 & -1 & 0 & -1\\
-1 & 3 & -1 & -1\\
0 & -1 & 2 & -1\\
-1 & -1 & -1 & 3
\end{bmatrix}.
\]

Notice the structure: degrees on the diagonal, $-1$ wherever there is an edge, and -- the
key property -- \textbf{every row sums to zero}. Of course it does: the diagonal is the
degree, which is exactly how many $-1$'s sit in that row. Zero row sums mean

\[
\mat{L}\,\vec{1} = \vec{0}, \qquad \vec{1} = (1,1,1,1)\T.
\]

So the all-ones vector is an eigenvector of $\mat{L}$ with eigenvalue $0$. \emph{Every}
Laplacian has $0$ as an eigenvalue. That single fact is the seed of everything that
follows.

\subsection{Why ``Laplacian''? The connection to calculus}

The name is borrowed from the Laplacian operator $\nabla^2 = \partial_{xx} + \partial_{yy}$
of multivariable calculus, and the borrowing is exact, not loose.

\begin{connectionbox}
On a grid of points spaced $h$ apart, the second derivative is approximated by
$f''(x) \approx \frac{f(x-h) - 2f(x) + f(x+h)}{h^2}$. That numerator -- ``twice the value
here, minus the value at each neighbor'' -- \emph{is} a row of the graph Laplacian of the
path graph, with the sign flipped. Acting on a vector of node-values,
\[
(\mat{L}\vec{x})_i = d_i\, x_i - \sum_{j \sim i} x_j = \sum_{j \sim i} (x_i - x_j),
\]
which compares the value at node $i$ to the average of its neighbors. The continuous
Laplacian measures how much a function at a point exceeds its surrounding average; the
graph Laplacian does the discrete same. Heat diffusion, $\partial_t u = -\mat{L}u$,
smooths a function on a graph exactly as the heat equation smooths temperature on a metal
plate. The whole field of \emph{spectral graph theory} is multivariable calculus done on
a graph instead of a continuum.
\end{connectionbox}

There is a second, cleaner way to see the Laplacian, and it explains why it is always
``nice'' (symmetric, never negative). Recall the incidence matrix $\mat{B}$ from
\S\ref{ch35:incidence}, whose rows are edge-differences. Form $\mat{B}\T \mat{B}$:

\begin{theorem}[The Laplacian factorization]
\label{ch35:thm-btb}
For any orientation of the edges, $\mat{L} = \mat{B}\T \mat{B}$.
\end{theorem}

\begin{proof}
The $(i,j)$ entry of $\mat{B}\T \mat{B}$ is the dot product of columns $i$ and $j$ of
$\mat{B}$ -- equivalently, $\sum_e B_{ei}B_{ej}$, summing over edges. Take the diagonal
entry $(i,i)$: $\sum_e B_{ei}^2$. Each edge touching node $i$ contributes a $\pm 1$ in that
column, so $B_{ei}^2 = 1$ for each of the $d_i$ edges at $i$ and $0$ otherwise; the sum is
$d_i$. Now take $i \ne j$: a term $B_{ei}B_{ej}$ is nonzero only for an edge $e$ touching
\emph{both} $i$ and $j$, i.e. the edge between them (if it exists). For that edge one
endpoint carries $+1$ and the other $-1$, so the product is $-1$. Hence the off-diagonal is
$-1$ for adjacent nodes and $0$ otherwise. That is exactly $\mat{L} = \mat{D}-\mat{A}$.
\end{proof}

Let us double-check on the running example by computing one entry. Column $2$ of $\mat{B}$
is $(+1,0,-1,-1,0)\T$ (edges $e_1,e_3,e_4$ touch node $2$). Its own dot product is
$1+1+1 = 3 = d_2$. Its dot product with column $1$, which is $(-1,-1,0,0,0)\T$, is
$(1)(-1) + 0 + 0 + 0 + 0 = -1$, matching $L_{12} = -1$. The factorization works.

\begin{aestheticbox}
The identity $\mat{L} = \mat{B}\T\mat{B}$ is small but it pays enormous dividends. Any
matrix of the form $\mat{M}\T\mat{M}$ is automatically \emph{symmetric} (so it has a full
set of orthogonal real eigenvectors) and \emph{positive semidefinite} (so its eigenvalues
are all $\ge 0$). Without computing a single eigenvalue we now know that the Laplacian's
spectrum is real, nonnegative, and orthogonally diagonalizable. The graph hands us a
perfectly behaved symmetric matrix for free. A great deal of applied mathematics works out
this way -- a messy object turns out to be $\mat{M}\T\mat{M}$ in disguise.
\end{aestheticbox}

\subsection{The energy form: what $\vec{x}\T \mat{L}\vec{x}$ measures}

Theorem~\ref{ch35:thm-btb} gives a one-line formula for the quadratic form. For any node
vector $\vec{x}$,
\[
\vec{x}\T \mat{L}\vec{x}
= \vec{x}\T \mat{B}\T \mat{B}\vec{x}
= (\mat{B}\vec{x})\T(\mat{B}\vec{x})
= \norm{\mat{B}\vec{x}}^2
= \sum_{(i,j)\in E} (x_i - x_j)^2.
\]

Read the last expression aloud: the Laplacian energy of an assignment $\vec{x}$ is the sum,
over all edges, of the squared difference between the two endpoints. It is small when
connected nodes get \emph{similar} values and large when an edge straddles very different
values.

\begin{focusbox}
$\vec{x}\T\mat{L}\vec{x} = \sum_{(i,j)\in E}(x_i - x_j)^2 \ge 0$ is the single formula to
remember about the Laplacian. It is ``how wiggly is $\vec{x}$ along the edges of the
graph.'' Minimizing it (subject to constraints) is the heart of spectral clustering,
smoothing, and embedding. Every Laplacian eigenvalue is a value of this energy on a
unit vector.
\end{focusbox}

Figure~\ref{ch35:fig-cut} makes the formula visible on our running graph. Label nodes
$\{1,2\}$ with $+1$ and nodes $\{3,4\}$ with $-1$. An edge inside a group contributes
$(x_i-x_j)^2 = 0$; an edge \emph{crossing} between the groups contributes
$(+1-(-1))^2 = 4$. Three edges cross, so the energy is $3 \times 4 = 12$ -- four times the
number of edges cut.

\begin{figure}[H]
\centering
\begin{tikzpicture}[scale=1.6,
  pnode/.style={circle,draw=sourceblue,very thick,fill=blue!10,minimum size=8mm,inner sep=0pt,font=\small},
  mnode/.style={circle,draw=trapred,very thick,fill=red!8,minimum size=8mm,inner sep=0pt,font=\small},
  inedge/.style={draw=spacegray,very thick},
  cutedge/.style={draw=trapred,very thick}]
  \node[pnode] (n1) at (0,1.3) {$1^{+}$};
  \node[pnode] (n2) at (1.5,1.3) {$2^{+}$};
  \node[mnode] (n3) at (1.5,0) {$3^{-}$};
  \node[mnode] (n4) at (0,0) {$4^{-}$};
  \draw[inedge] (n1) -- (n2);
  \draw[cutedge] (n1) -- (n4);
  \draw[cutedge] (n2) -- (n3);
  \draw[cutedge] (n2) -- (n4);
  \draw[inedge] (n3) -- (n4);
  % cut line
  \draw[dashed,trapred,thick] (-0.55,0.65) -- (2.05,0.65);
  \node[font=\scriptsize,trapred] at (2.45,0.65) {cut};
\end{tikzpicture}
\caption{The energy $\vec{x}\T\mat{L}\vec{x}$ as a cut counter. With $\vec{x}=(+1,+1,-1,-1)\T$,
the two gray edges lie inside a group and cost $0$; the three red edges cross the dashed cut
and each cost $(+1-(-1))^2 = 4$. Total energy $= 12 = 4 \times (\text{edges cut})$. Low energy
means few edges straddle the boundary -- the link between the Laplacian and clustering.}
\label{ch35:fig-cut}
\end{figure}

Now we can see \emph{why} $0$ is the smallest eigenvalue and what makes it special. The
energy is zero exactly when $x_i = x_j$ across every edge -- i.e. $\vec{x}$ is constant on
each connected piece of the graph. We met this in \S\ref{ch35:incidence}: it is the
nullspace of $\mat{B}$, hence of $\mat{L}$.

\section{Eigenvalues of the Laplacian: hearing connectivity}
\label{ch35:spectrum}

Because $\mat{L}$ is symmetric positive semidefinite, we can order its eigenvalues:
\[
0 = \lambda_1 \le \lambda_2 \le \cdots \le \lambda_n.
\]
We always have $\lambda_1 = 0$ (the constant eigenvector). The next theorem is the jewel of
the subject: the \emph{number of zero eigenvalues counts the connected components}.

\begin{theorem}[Multiplicity of zero $=$ number of components]
\label{ch35:thm-components}
The multiplicity of the eigenvalue $0$ of the graph Laplacian equals the number of
connected components of the graph.
\end{theorem}

\begin{proof}
An eigenvector for $0$ is a vector $\vec{x}$ with $\mat{L}\vec{x}=\vec{0}$, equivalently
$\vec{x}\T\mat{L}\vec{x} = \sum_{(i,j)\in E}(x_i-x_j)^2 = 0$ (using that $\mat{L}$ is
positive semidefinite, energy zero forces $\mat{L}\vec x = 0$). Every squared term must
vanish, so $x_i = x_j$ across each edge. Hence $\vec{x}$ is constant on each connected
component, and free to choose a different constant per component. If there are $c$
components, the space of such vectors has dimension exactly $c$ -- one indicator vector per
component, and these are linearly independent. So the eigenspace of $0$ has dimension $c$.
\end{proof}

\begin{intuitionbox}
Here is the picture. A connected graph holds together with one ``ground state'' (everybody
the same), so $0$ appears once. Break the graph into two islands and you get \emph{two}
independent ground states -- color island A red, island B blue, and there is no edge to
penalize the jump. Each island contributes its own zero. The eigenvalue $0$ is literally
counting the islands.
\end{intuitionbox}

This already tells us something operational: to test whether a network is connected, you do
not trace paths -- you compute the second-smallest eigenvalue $\lambda_2$. If
$\lambda_2 > 0$, the graph is connected. If $\lambda_2 = 0$, it has fallen apart.

\subsection{The Fiedler value: connectivity by degree}

The eigenvalue $\lambda_2$ is so important it has a name: the \textbf{algebraic
connectivity}, or \textbf{Fiedler value}, after Miroslav Fiedler. It does not just say
\emph{whether} the graph is connected; its \emph{size} measures \emph{how robustly}
connected the graph is. A large $\lambda_2$ means the graph is hard to cut into pieces; a
small positive $\lambda_2$ means there is a near-bottleneck -- a thin bridge that almost
disconnects it.

\begin{figure}[H]
\centering
\begin{tikzpicture}[scale=1.1,
  node/.style={circle,draw=sourceblue,thick,fill=blue!8,minimum size=6mm,inner sep=0pt,font=\scriptsize},
  every edge/.style={draw=spacegray,thick}]
  % left cluster
  \node[node] (a1) at (1,1) {};
  \node[node] (a2) at (0,1.6) {};
  \node[node] (a3) at (0,0.4) {};
  \draw (a1) edge (a2); \draw (a1) edge (a3); \draw (a2) edge (a3);
  % right cluster
  \node[node] (b1) at (3,1) {};
  \node[node] (b2) at (4,1.6) {};
  \node[node] (b3) at (4,0.4) {};
  \draw (b1) edge (b2); \draw (b1) edge (b3); \draw (b2) edge (b3);
  % bridge
  \draw[trapred,very thick,dashed] (a1.east) -- (b1.west) ;
  \node[font=\scriptsize,trapred] at (2,1.25) {bridge};
\end{tikzpicture}
\caption{Two tight triangles joined by a single ``bridge'' edge. The graph is connected, so
$\lambda_2 > 0$, but only barely: cutting the one dashed edge splits it. The Fiedler value
$\lambda_2$ is small, and the Fiedler eigenvector takes one sign on the left cluster and the
opposite sign on the right -- it \emph{finds} the cut.}
\label{ch35:fig-bridge}
\end{figure}

Why does the eigenvector know where to cut? Because of the energy interpretation. The
Fiedler eigenvector $\vec{v}_2$ minimizes $\vec{x}\T\mat{L}\vec{x} = \sum (x_i-x_j)^2$ over
all unit vectors \emph{orthogonal to the constant vector} $\vec{1}$ (orthogonality to the
$\lambda_1$-eigenvector is forced by the spectral theorem). ``Orthogonal to $\vec{1}$''
means $\sum_i x_i = 0$: the values must be partly positive and partly negative. So $\vec{v}_2$
is the gentlest possible nonconstant labeling -- it wants neighbors to share a value, but it
is forced to use both signs. The cheapest way to do that is to make the jump happen across
\emph{few} edges, i.e. across a bottleneck. Splitting nodes by the \emph{sign} of $\vec{v}_2$
is the classic \textbf{spectral bisection} heuristic.

\begin{labbox}
Take the path graph on four nodes in a line, $1\!-\!2\!-\!3\!-\!4$. Its Laplacian is
\[
\mat{L} =
\begin{bmatrix} 1&-1&0&0\\ -1&2&-1&0\\ 0&-1&2&-1\\ 0&0&-1&1 \end{bmatrix}.
\]
Verify $\mat{L}\vec{1}=\vec0$. Then test the simple sign pattern $\vec{s} = (1,1,-1,-1)\T$,
which splits the path in half: compute
\[
\mat{L}\,\vec{s} = (0,\,2,\,-2,\,0)\T.
\]
This is \emph{not} a scalar multiple of $\vec{s}$, so $\vec{s}$ is not exactly an
eigenvector -- but notice it already carries the right idea, putting all the ``action'' at
the two middle nodes where the cut happens. The true Fiedler vector of the $4$-path is
$\vec{v}_2 = \frac{1}{\sqrt{2}}(\cos\tfrac{\pi}{8},\,\cos\tfrac{3\pi}{8},\,-\cos\tfrac{3\pi}{8},\,
-\cos\tfrac{\pi}{8})\T \approx (0.65,\,0.27,\,-0.27,\,-0.65)\T$: its values decrease smoothly
from one end to the other, and its sign change happens in the middle, telling you to cut the
path in half. Its eigenvalue is $\lambda_2 = 2 - 2\cos\tfrac{\pi}{4} = 2 - \sqrt2 \approx
0.586$. (The exact eigenvalues of the $n$-node path are $2 - 2\cos(\pi k/n)$ for
$k = 0,1,\dots,n-1$, a fact worth memorizing: it is the discrete cosine transform hiding in
a graph.)
\end{labbox}

\begin{applicationbox}
\textbf{Spectral clustering in machine learning.} Suppose you have a pile of data points
-- images, customers, documents -- and a notion of similarity between pairs. Build a graph:
nodes are data points, and you connect similar points with edges (often weighted, see
below). The communities in your data are the loosely-connected clusters of this graph.
Spectral clustering computes the bottom few Laplacian eigenvectors $\vec{v}_2,\dots,\vec{v}_k$,
uses them as new coordinates for each point (the $i$-th point becomes the row
$(\,v_2(i),\dots,v_k(i)\,)$), and runs plain $k$-means in that low-dimensional space. The
Laplacian eigenvectors \emph{unfold} the graph so that clusters that were tangled in the
original feature space become linearly separable blobs. This is the algorithm behind a huge
amount of image segmentation, community detection in social networks, and unsupervised
grouping. Its entire justification is Theorem~\ref{ch35:thm-components} and the energy
formula: small eigenvalues correspond to ``almost components,'' i.e. clusters.
\end{applicationbox}

\subsection{Weighted graphs}

Real networks rarely have all edges equal. A friendship can be strong or weak; a road can
be wide or narrow; a similarity can be near or far. Attach a weight $w_{ij} > 0$ to each
edge. Then the weighted adjacency matrix has $A_{ij} = w_{ij}$, the degree becomes the
\emph{weighted} degree $d_i = \sum_j w_{ij}$, and the Laplacian $\mat{L} = \mat{D}-\mat{A}$
keeps every property we proved -- because the energy form is now
$\vec{x}\T\mat{L}\vec{x} = \sum_{(i,j)\in E} w_{ij}(x_i - x_j)^2$, still a sum of
nonnegative terms. Heavy edges penalize differences more, so the eigenvectors try harder to
keep strongly-connected nodes together. Everything in this section runs unchanged with
weights; we will not belabor it, but it is what makes the machinery useful on real data.

\section{Random walks and Markov chains}
\label{ch35:markov}

We now change the question. Instead of a fixed labeling of nodes, imagine a \emph{walker}
hopping around the graph at random, and ask: where is the walker likely to be? This is the
setting of \textbf{Markov chains}, where transition matrices and the dominant eigenvector
take center stage.

\begin{intuitionbox}
A Markov chain is a memoryless random walk. The walker sits at some node, looks at the
outgoing edges, and picks one at random according to fixed probabilities -- then forgets
how it got there. ``Memoryless'' (the \emph{Markov property}) means the next step depends
only on \emph{where you are}, never on the history of how you arrived. The entire future is
encoded in one matrix of one-step probabilities.
\end{intuitionbox}

\begin{historybox}
In the early 1900s in Imperial Russia, an intense dispute erupted between two mathematicians. Pavel Nekrasov published papers claiming that statistical independence was an absolute necessity for the Law of Large Numbers to hold, arguing that independent probabilities reflected moral free will. Andrey Markov was appalled by what he saw as theological doctrine masquerading as mathematics. To refute Nekrasov, Markov set out in 1906 to prove that even strongly dependent sequences of events converge to predictable long-term averages. To prove it empirically, Markov spent months counting vowel-consonant transitions across the first 20,000 letters of Pushkin's poem \emph{Eugene Onegin}. He showed that a vowel followed a consonant with probability $0.872$ and another vowel with probability $0.128$, yet the overall vowel frequency stabilized rigorously at $43.2\%$---the birth of \textbf{Markov chains}.

Around the same time, Oskar Perron (1907) and Georg Frobenius (1912) proved the Perron--Frobenius theorem: any strictly positive matrix has a unique, positive leading eigenvalue with a strictly positive eigenvector. For row-stochastic transition matrices, this eigenvalue is always $\lambda_1 = 1$, guaranteeing a unique stationary probability distribution.

In 1998, Stanford graduate students Larry Page and Sergey Brin modeled the World Wide Web as an enormous Markov chain. Naive random walks on hyperlinks got trapped in dead ends or infinite loops (spider traps). Brin and Page solved this by adding a small teleportation probability ($1-\alpha \approx 0.15$) to jump to any page uniformly. This made the Google matrix $\mat{G} = \alpha \mat{P} + \frac{1-\alpha}{n} \vec{1}\vec{1}\T$ strictly positive, ensuring that power iteration converges reliably to a unique ranking vector of web page importance.
\end{historybox}

Make this precise with a small \emph{directed, weighted} example -- a tiny model of weather,
say, or of a user browsing three pages.

\begin{figure}[H]
\centering
\begin{tikzpicture}[scale=1.05,>=Stealth]
  % Panel 1: Transition probabilities
  \begin{scope}[shift={(0,0)}]
    \node[anchor=south,font=\footnotesize\scshape,text=spacegray!90!black] at (1.5,1.5) {1. Transition Probabilities $P_{ij}$};
    \draw[fill=codebg,draw=spacegray!30,rounded corners=4pt] (-1.4,-2.4) rectangle (4.4,1.4);

    \node[circle,draw=sourceblue,very thick,fill=sourceblue!12,minimum size=8mm,inner sep=0pt,font=\small] (s) at (0,0) {$S$};
    \node[circle,draw=targetgreen,very thick,fill=targetgreen!12,minimum size=8mm,inner sep=0pt,font=\small] (c) at (3.0,0) {$C$};
    \node[circle,draw=quotientorange,very thick,fill=quotientorange!12,minimum size=8mm,inner sep=0pt,font=\small] (r) at (1.5,-1.7) {$R$};

    \draw[->,thick,spacegray!80!black] (s) to[bend left=18] node[midway,above=1pt,font=\scriptsize,fill=codebg,inner sep=1pt]{$0.5$} (c);
    \draw[->,thick,spacegray!80!black] (c) to[bend left=18] node[midway,below=1pt,font=\scriptsize,fill=codebg,inner sep=1pt]{$0.3$} (s);
    \draw[->,thick,spacegray!80!black] (s) to[bend left=12] node[midway,left=2pt,font=\scriptsize,fill=codebg,inner sep=1pt]{$0.3$} (r);
    \draw[->,thick,spacegray!80!black] (r) to[bend left=12] node[midway,right=1pt,font=\scriptsize,fill=codebg,inner sep=1pt]{$0.4$} (c);
    \draw[->,thick,spacegray!80!black] (c) to[bend left=12] node[midway,right=2pt,font=\scriptsize,fill=codebg,inner sep=1pt]{$0.4$} (r);
    \draw[->,thick,spacegray!80!black] (r) to[bend left=12] node[midway,left=1pt,font=\scriptsize,fill=codebg,inner sep=1pt]{$0.4$} (s);

    \draw[->,thick,spacegray!80!black] (s) edge[loop left,looseness=4] node[midway,left=1pt,font=\scriptsize]{$0.2$} (s);
    \draw[->,thick,spacegray!80!black] (c) edge[loop right,looseness=4] node[midway,right=1pt,font=\scriptsize]{$0.3$} (c);
    \draw[->,thick,spacegray!80!black] (r) edge[loop below,looseness=4] node[midway,below=1pt,font=\scriptsize]{$0.2$} (r);
  \end{scope}

  % Panel 2: Steady state probability fluid flow
  \begin{scope}[shift={(6.8,0)}]
    \node[anchor=south,font=\footnotesize\scshape,text=spacegray!90!black] at (1.5,1.5) {2. Steady-State Flux Equilibrium ($\vec{\pi}^*\mat{P} = \vec{\pi}^*$)};
    \draw[fill=codebg,draw=spacegray!30,rounded corners=4pt] (-1.2,-2.4) rectangle (4.2,1.4);

    % Nodes sized and labeled with stationary probabilities pi* = (0.301, 0.391, 0.308)
    \node[circle,draw=sourceblue,line width=1.5pt,fill=sourceblue!15,minimum size=9.5mm,inner sep=0pt,font=\footnotesize,align=center] (s2) at (0,0) {$S$\\[-2pt]\tiny $30.1\%$};
    \node[circle,draw=targetgreen,line width=1.8pt,fill=targetgreen!15,minimum size=11.5mm,inner sep=0pt,font=\footnotesize,align=center] (c2) at (3.0,0) {$C$\\[-2pt]\tiny $39.1\%$};
    \node[circle,draw=quotientorange,line width=1.5pt,fill=quotientorange!15,minimum size=9.8mm,inner sep=0pt,font=\footnotesize,align=center] (r2) at (1.5,-1.7) {$R$\\[-2pt]\tiny $30.8\%$};

    % Directed flux corridors with net exchange
    \draw[->,line width=1.4pt,vecblue] (s2) to[bend left=16] node[midway,above,font=\tiny,text=vecblue]{$15.0\%$} (c2);
    \draw[->,line width=1.2pt,targetgreen] (c2) to[bend left=16] node[midway,below,font=\tiny,text=targetgreen]{$11.7\%$} (s2);
    \draw[->,line width=1.2pt,quotientorange] (r2) to[bend left=12] node[midway,left,font=\tiny,text=quotientorange]{$12.3\%$} (s2);
    \draw[->,line width=1.4pt,targetgreen] (c2) to[bend left=12] node[midway,right,font=\tiny,text=targetgreen]{$15.6\%$} (r2);

    \node[draw=spacegray!40,fill=white,rounded corners=2pt,font=\tiny,align=center] at (1.5,-0.75) {Inflow $=$ Outflow\\at every node};
  \end{scope}
\end{tikzpicture}
\caption{A three-state weather Markov chain and its steady-state probability equilibrium. Left: Transition probabilities $P_{ij}$ for Sunny ($S$), Cloudy ($C$), Rainy ($R$), where every node's outgoing arrows sum to $1$. Right: The stationary distribution $\vec{\pi}^* \approx (0.301, 0.391, 0.308)$ acts as an invariant probability fluid. At equilibrium, the probability mass entering any state via incoming flows exactly balances the probability mass exiting via outgoing flows, making $\vec{\pi}^*$ the left eigenvector for $\lambda=1$.}
\label{ch35:fig-markov}
\end{figure}

Collect the transition probabilities into a matrix. We will use the \textbf{row-stochastic}
convention: $P_{ij}$ is the probability of going \emph{from} state $i$ \emph{to} state $j$,
so each \emph{row} sums to $1$. From Figure~\ref{ch35:fig-markov}, ordering states
$S,C,R$:

\[
\mat{P} =
\begin{bmatrix}
P_{SS} & P_{SC} & P_{SR}\\
P_{CS} & P_{CC} & P_{CR}\\
P_{RS} & P_{RC} & P_{RR}
\end{bmatrix}
=
\begin{bmatrix}
0.2 & 0.5 & 0.3\\
0.3 & 0.3 & 0.4\\
0.4 & 0.4 & 0.2
\end{bmatrix}.
\]

Check the rows: $0.2+0.5+0.3 = 1$, $0.3+0.3+0.4=1$, $0.4+0.4+0.2=1$. A matrix with
nonnegative entries and rows summing to $1$ is called a \textbf{stochastic matrix} (or
\emph{transition matrix}). Its rows are probability distributions.

\begin{focusbox}
A stochastic matrix has nonnegative entries and \emph{rows summing to $1$}. The all-ones
vector is a \emph{right} eigenvector with eigenvalue $1$: $\mat{P}\vec{1} = \vec{1}$, because
each row sums to $1$. Hold this thought -- the interesting eigenvector, the stationary
distribution, will be the \emph{left} eigenvector for the same eigenvalue $1$.
\end{focusbox}

\subsection{Distributions evolve by matrix multiplication}

Represent the walker's location as a \emph{row} probability vector
$\vec{\pi} = (\pi_S, \pi_C, \pi_R)$ with $\pi_i \ge 0$ and $\sum_i \pi_i = 1$. Say today we
are certain it is sunny: $\vec{\pi}^{(0)} = (1,0,0)$. What is the distribution tomorrow?
The probability of being cloudy tomorrow is the probability we are in each state today times
the probability that state goes to cloudy -- a dot product with column $C$ of $\mat{P}$.
Stacked together, \emph{the update is a vector--matrix product}:

\[
\vec{\pi}^{(1)} = \vec{\pi}^{(0)}\mat{P}.
\]

With $\vec{\pi}^{(0)} = (1,0,0)$ this just reads off the first row:
$\vec{\pi}^{(1)} = (0.2, 0.5, 0.3)$. After two days,
$\vec{\pi}^{(2)} = \vec{\pi}^{(1)}\mat{P}$:

\[
\vec{\pi}^{(2)} = (0.2,\,0.5,\,0.3)
\begin{bmatrix}
0.2 & 0.5 & 0.3\\
0.3 & 0.3 & 0.4\\
0.4 & 0.4 & 0.2
\end{bmatrix}
= (0.31,\ 0.37,\ 0.32).
\]

Let us check the first entry: $0.2(0.2) + 0.5(0.3) + 0.3(0.4) = 0.04 + 0.15 + 0.12 = 0.31$.
Good -- and notice the new vector still sums to $1$ ($0.31+0.37+0.32 = 1$), as it must: a
stochastic matrix maps probability distributions to probability distributions. In general
$\vec{\pi}^{(k)} = \vec{\pi}^{(0)}\mat{P}^k$. Evolving the chain $k$ steps is multiplying by
$\mat{P}$ a total of $k$ times.

\begin{trapbox}
Mind the convention. With the \emph{row-stochastic} matrix above, distributions are
\emph{row} vectors and you multiply on the \emph{left}: $\vec\pi \mapsto \vec\pi\mat{P}$.
Many texts (and most physics writing) instead use \emph{column}-stochastic matrices, where
columns sum to $1$, distributions are column vectors, and you multiply on the right:
$\vec\pi \mapsto \mat{P}\vec\pi$. The two are transposes of each other. Pick one and stay
with it for a whole calculation; mixing them silently transposes your transition
probabilities and produces nonsense that still sums to $1$, so it is a maddening bug to
catch. In this chapter: rows sum to $1$, distributions on the left.
\end{trapbox}

\subsection{The stationary distribution is a left eigenvector}

Run the chain a few more steps and watch:

\[
\begin{aligned}
\vec\pi^{(2)} &= (0.310,\ 0.370,\ 0.320),\\
\vec\pi^{(3)} &= (0.301,\ 0.394,\ 0.305),\\
\vec\pi^{(4)} &= (0.300,\ 0.391,\ 0.309),\\
&\ \,\vdots
\end{aligned}
\]

The numbers are settling down. Figure~\ref{ch35:fig-converge} plots each coordinate against
the step number: starting from the spike ``certainly sunny,'' all three probabilities swing
once and then lock onto fixed levels.

\begin{figure}[H]
\centering
\begin{tikzpicture}
\begin{axis}[
  width=11cm, height=6cm,
  title={Convergence of State Probabilities $\vec{\pi}^{(k)} \to \vec{\pi}^*$},
  title style={font=\small},
  xlabel={step $k$}, ylabel={probability $\pi^{(k)}$},
  xmin=0, xmax=7, ymin=0, ymax=1.05,
  xtick={0,1,2,3,4,5,6,7},
  legend pos=north east, legend cell align=left,
  grid=major, grid style={gray!20},
  every axis plot/.append style={very thick,mark=*,mark size=1.4pt}]
  \addplot[sourceblue] coordinates
    {(0,1.0000)(1,0.2000)(2,0.3100)(3,0.3010)(4,0.3004)(5,0.3009)(6,0.3007)(7,0.3008)};
  \addlegendentry{$\pi_S$ (sunny)}
  \addplot[targetgreen] coordinates
    {(0,0.0000)(1,0.5000)(2,0.3700)(3,0.3940)(4,0.3907)(5,0.3910)(6,0.3910)(7,0.3910)};
  \addlegendentry{$\pi_C$ (cloudy)}
  \addplot[trapred] coordinates
    {(0,0.0000)(1,0.3000)(2,0.3200)(3,0.3050)(4,0.3089)(5,0.3082)(6,0.3083)(7,0.3083)};
  \addlegendentry{$\pi_R$ (rainy)}
  \draw[gray!60,dashed] (axis cs:0,0.3910) -- (axis cs:7,0.3910);
\end{axis}
\end{tikzpicture}
\caption{The weather chain forgets its start. Each coordinate of $\vec\pi^{(k)}$ converges to
the stationary value $\vec\pi^\ast \approx (0.301,\,0.391,\,0.308)$ within a handful of steps,
no matter that we began ``certainly sunny.'' The geometric rate of approach is set by the
second eigenvalue $|\lambda_2|$ of $\mat{P}$.}
\label{ch35:fig-converge}
\end{figure}

They are converging to a vector $\vec\pi^\ast$ that does not
change when you push it through $\mat{P}$ one more time:

\[
\vec\pi^\ast \mat{P} = \vec\pi^\ast.
\]

This is the \textbf{stationary distribution}. Algebraically it says $\vec\pi^\ast$ is a
\emph{left} eigenvector of $\mat{P}$ with eigenvalue $1$ -- equivalently a right eigenvector
of $\mat{P}\T$. There is the headline of this section, stated as plainly as it can be:

\begin{focusbox}
\textbf{The long-run behavior of a Markov chain is an eigenvector problem.} The fraction of
time the walker spends in each state, after wandering forever, is the left eigenvector of
the transition matrix $\mat{P}$ for eigenvalue $1$, normalized to sum to $1$. Find that
eigenvector and you have predicted the future of the whole system.
\end{focusbox}

Let us actually solve $\vec\pi^\ast\mat{P} = \vec\pi^\ast$ for our weather chain. Writing
$\vec\pi^\ast = (a,b,c)$, the equation $\vec\pi^\ast(\mat{P}-\mat{I}) = \vec0$ together with
$a+b+c=1$ gives the system (one equation per column of $\mat P$, dropping one as
redundant):
\[
\begin{aligned}
0.2a + 0.3b + 0.4c &= a,\\
0.5a + 0.3b + 0.4c &= b,\\
a + b + c &= 1.
\end{aligned}
\]
From the first equation, $0.3b + 0.4c = 0.8a$. From the second,
$0.5a + 0.4c = 0.7b$. Solving this $3\times3$ system (a good by-hand exercise) yields the
exact answer
\[
\vec\pi^\ast = \left(\tfrac{40}{133},\ \tfrac{52}{133},\ \tfrac{41}{133}\right)
\approx (0.3008,\ 0.3910,\ 0.3083),
\]
exactly the vector the iteration was crawling toward. In the long run the weather is sunny
about $30\%$ of days, cloudy about $39\%$, rainy about $31\%$ -- regardless of today's
weather. The starting vector washed out completely. That ``forgetting of the start'' is the
defining feature of a well-behaved chain.

\subsection{When is there a unique stationary distribution it converges to?}

Not every chain settles down to a single answer. Two things can go wrong, and a beautiful
theorem says they are the only two.

\begin{formalbox}
\textbf{Perron--Frobenius for stochastic matrices.} If $\mat{P}$ is the transition matrix of
a finite Markov chain that is \emph{irreducible} (you can get from any state to any other,
following arrows) and \emph{aperiodic} (it is not locked into a fixed cycle of return
times), then:
\begin{enumerate}
\item $\mat{P}$ has a unique stationary distribution $\vec\pi^\ast$, a positive left
eigenvector for the eigenvalue $1$;
\item $1$ is the largest eigenvalue in magnitude, and it is simple (multiplicity one);
\item for \emph{every} starting distribution, $\vec\pi^{(0)}\mat{P}^k \to \vec\pi^\ast$ as
$k \to \infty$.
\end{enumerate}
The rate of convergence is governed by the \emph{second-largest} eigenvalue $\lambda_2$ in
magnitude: the distance to $\vec\pi^\ast$ shrinks like $|\lambda_2|^k$. This gap
$1 - |\lambda_2|$ is the \emph{spectral gap}, the mixing speed of the chain.
\end{formalbox}

\begin{intuitionbox}
\emph{Reducible} means the graph splits into pieces a walker can get stuck inside -- then
the long-run answer depends on where you started, so there is no single stationary
distribution. \emph{Periodic} means the chain marches in lockstep (think bipartite: odd
steps here, even steps there); the distribution oscillates forever instead of settling.
Rule out both -- by making sure every state can reach every other and that there is at least
one self-loop or odd cycle to break the rhythm -- and the chain is guaranteed to converge to
one universal distribution. The weather chain has all positive transitions, so it is
strictly irreducible and aperiodic, explaining why it converged so cleanly.
\end{intuitionbox}

\begin{connectionbox}
The eigenvalue $\lambda_2$ controlling Markov mixing and the eigenvalue $\lambda_2$
controlling graph connectivity (the Fiedler value) are \emph{the same idea in two
costumes}. For a random walk on an undirected graph, the transition matrix is
$\mat{P} = \mat{D}\inv\mat{A}$, and its eigenvalues are related to the (normalized)
Laplacian eigenvalues by $\mu_k = 1 - \lambda_k$. A large Laplacian gap means a large
Markov gap means fast mixing means no bottleneck. A graph that is hard to cut is a graph a
random walker mixes through quickly. Connectivity, clustering, and mixing time are one
spectral story told from three angles, and that unity is what gives spectral methods their
reach.
\end{connectionbox}

\section{PageRank: importance as an eigenvector}
\label{ch35:pagerank}

We close with the application that made all of this famous: Google's PageRank. The problem
is to rank the pages of the web by ``importance'' using only the link structure -- no human
judgment, no reading the text, just who links to whom.

\begin{intuitionbox}
The founding idea is delightfully circular and entirely correct: \emph{a page is important
if important pages link to it}. A link is a vote, but votes are weighted by the importance
of the voter, and the voter's importance is decided the same way. That circular definition
is not a paradox -- it is an eigenvector equation. ``$\vec{x}$ is proportional to the
incoming weighted votes, which are proportional to $\vec{x}$'' is precisely
$\mat{M}\vec{x} = \vec{x}$.
\end{intuitionbox}

Model the web as a directed graph: nodes are pages, an edge $i \to j$ means page $i$ links
to page $j$. Imagine a \textbf{random surfer} who, at each step, picks uniformly at random
one of the outgoing links of the current page and clicks it. If page $i$ has $d_i$ outgoing
links, the transition probability from $i$ to each linked page $j$ is $1/d_i$. This is just
the Markov chain $\mat{P}$ from the previous section, built from the link graph. PageRank is
\emph{the stationary distribution of the random surfer}: the fraction of time an
endlessly-clicking surfer spends on each page. Important pages are the ones the surfer keeps
landing on.

Take a tiny four-page web. Page $1$ links to $2$ and $3$; page $2$ links to $3$; page $3$
links to $1$; page $4$ links to $3$.

\begin{figure}[H]
\centering
\begin{tikzpicture}[scale=1.7,
  node/.style={circle,draw=sourceblue,very thick,fill=blue!8,minimum size=9mm,inner sep=0pt,font=\small},
  every edge/.style={draw=spacegray,thick,-{Stealth[length=2mm]}}]
  \node[node] (p1) at (0,1.2) {$1$};
  \node[node] (p2) at (1.5,1.2) {$2$};
  \node[node] (p3) at (1.5,0) {$3$};
  \node[node] (p4) at (0,0) {$4$};
  \draw (p1) edge (p2);
  \draw (p1) edge (p3);
  \draw (p2) edge (p3);
  \draw (p3) edge[bend left=15] (p1);
  \draw (p4) edge (p3);
\end{tikzpicture}
\caption{A four-page web. Page $3$ is linked from $1$, $2$, and $4$ -- it should rank
highest. Page $1$ gets its importance from page $3$, which is itself highly ranked, so $1$
does well too despite only one inbound link. PageRank captures exactly this ``importance
flows backward along links'' intuition.}
\label{ch35:fig-web}
\end{figure}

Page $1$ has two out-links, so it sends probability $1/2$ to each of $2$ and $3$. Pages $2$,
$3$, $4$ each have a single out-link and send all their probability along it. The
row-stochastic transition matrix is

\[
\mat{P} =
\begin{bmatrix}
0 & \tfrac12 & \tfrac12 & 0\\[2pt]
0 & 0 & 1 & 0\\[2pt]
1 & 0 & 0 & 0\\[2pt]
0 & 0 & 1 & 0
\end{bmatrix}.
\]

We want the left eigenvector for eigenvalue $1$. But there is a problem lurking, and the fix
for it is the cleverest part of the whole algorithm.

\subsection{Two snags and the ``teleport'' that fixes both}

\textbf{Dangling nodes.} A page with \emph{no} outgoing links (a PDF, a dead end) gives a
zero row -- the surfer falls off the edge of the web and probability leaks away; $\mat{P}$
is no longer stochastic. \textbf{Reducibility and periodicity.} Even with no dead ends, the
real web has spider-traps and disconnected islands, so Perron--Frobenius does not apply and
there may be no unique stationary distribution.

Brin and Page fixed both at once with a single idea: \textbf{teleportation}. With
probability $\alpha$ (the \emph{damping factor}, classically $\alpha = 0.85$) the surfer
follows a random link as before; with probability $1-\alpha$ the surfer gets bored and jumps
to a page chosen \emph{uniformly at random} from all $n$ pages. The new transition matrix is

\[
\mat{G} = \alpha \mat{P} + (1-\alpha)\,\tfrac1n\,\mat{J},
\]

where $\mat{J}$ is the all-ones $n\times n$ matrix (and dangling rows of $\mat P$ are first
replaced by uniform rows $\tfrac1n\vec1\T$, so the surfer ``teleports out'' of dead ends).
This $\mat{G}$ is the \textbf{Google matrix}.

\begin{aestheticbox}
Look at what teleportation buys, all from one convex combination. Because $\mat{J}/n$ has
\emph{every} entry positive, $\mat{G}$ has every entry strictly positive: from any page you
can reach any other in one step (you might teleport there). Strict positivity makes
$\mat{G}$ irreducible \emph{and} aperiodic in a single stroke, so Perron--Frobenius applies
with no exceptions: a unique, positive, stationary PageRank vector \emph{always} exists and
the iteration \emph{always} converges. One small probabilistic tweak turns a fragile,
ill-defined problem into one with a guaranteed clean answer.
\end{aestheticbox}

The PageRank vector $\vec{r}$ is the stationary distribution of $\mat{G}$:
\[
\vec{r}\,\mat{G} = \vec{r}, \qquad \sum_i r_i = 1.
\]
Equivalently, written as the update each entry must satisfy (this is the form you will see
in code),
\[
r_j = \frac{1-\alpha}{n} + \alpha \sum_{i \to j} \frac{r_i}{d_i},
\]
``my rank is a baseline teleport share plus a weighted sum of the ranks flowing in from
pages that link to me, each split evenly among its out-links.'' Read it twice; it is the
whole algorithm.

\subsection{The power method: how it is actually computed}

The web has billions of pages; you cannot diagonalize a billion-by-billion matrix. But you
do not need to. You only need the \emph{dominant} eigenvector, and there is a one-line
method for that, which you have already been using without naming it: just iterate.

\begin{intuitionbox}
\textbf{The power method.} Start from any distribution $\vec\pi^{(0)}$ (say uniform) and
repeatedly multiply: $\vec\pi^{(k+1)} = \vec\pi^{(k)}\mat{G}$. Write the start in the
eigenvector basis; each multiplication multiplies the component along eigenvector
$\vec{v}_k$ by $\lambda_k$. Since $\lambda_1 = 1$ dominates and every other $|\lambda_k| < 1$
shrinks toward zero, all the other components die off and what survives is the
$\lambda_1$-eigenvector -- the PageRank vector. Repeated multiplication is an
eigenvector-finding machine that automatically keeps the biggest eigenvalue and forgets the
rest.
\end{intuitionbox}

This is exactly the ``run the chain forward and watch it settle'' computation from
\S\ref{ch35:markov}, now recognized as the power method. Each step is one sparse
matrix--vector product -- cheap, because each page has only a handful of links. Convergence
is geometric at rate $\alpha = 0.85$ (the spectral gap of $\mat{G}$ is controlled by the
damping factor; in fact $|\lambda_2| \le \alpha$ for the Google matrix), so a few dozen
iterations suffice to rank the entire web. That is why the choice $\alpha = 0.85$ is a
tuning knob: lower $\alpha$ means faster convergence but more weight on random teleporting
and less on the actual link structure.

\begin{labbox}
Run three steps of the power method by hand on the four-page web of
Figure~\ref{ch35:fig-web}, using the \emph{undamped} $\mat{P}$ for simplicity and starting
from uniform $\vec\pi^{(0)} = (\tfrac14,\tfrac14,\tfrac14,\tfrac14)$. Compute
$\vec\pi^{(1)} = \vec\pi^{(0)}\mat{P}$, then $\vec\pi^{(2)}$, then $\vec\pi^{(3)}$. Watch
page $3$'s entry climb and page $4$'s entry collapse toward $0$ (nothing links to page $4$,
so the surfer never returns to it). You are watching importance flow through the graph in
real time. Then reflect: with damping, page $4$ would keep a small floor of
$\tfrac{1-\alpha}{n}$ -- teleportation guarantees no page is ranked exactly zero.
\end{labbox}

\begin{applicationbox}
\textbf{PageRank is everywhere now, not just search.} The same ``stationary distribution of
a smartly-chosen random walk'' template reappears across modern computing:
\emph{personalized PageRank} (teleport back to one user's pages, not uniformly) powers
friend and product recommendations; \emph{TextRank} runs PageRank on a graph of sentences to
do automatic summarization; \emph{citation analysis} ranks scientific papers and journals;
biologists rank genes and proteins in interaction networks (``GeneRank''); and random-walk
node embeddings (DeepWalk, node2vec) -- the graph analog of word embeddings -- feed graph
neural networks. Every one of them is the chapter you just read: build a transition matrix
from a network, find the eigenvector for eigenvalue $1$. The eigenvector of a giant sparse
matrix turned out to be one of the most valuable objects of the information age.
\end{applicationbox}

\begin{walkbox}
Take this on a walk. You have now seen the eigenvalue $1$ play the starring role twice -- as
the stationary distribution of a Markov chain, and as PageRank -- and the eigenvalue $0$
play it twice as well -- counting connected components, and the constant ground state of the
Laplacian. Why is it that ``the boring eigenvalue'' ($0$ for $\mat{L}$, $1$ for $\mat{P}$)
carries the most important information? What is the relationship between $\mat{L} =
\mat{D}-\mat{A}$ and $\mat{P} = \mat{D}\inv\mat{A}$, and how does the eigenvalue $0$ of one
become the eigenvalue $1$ of the other? (Hint: $\mat{P} = \mat{I} - \mat{D}\inv\mat{L}$.)
Hold that identity in your head and the two halves of this chapter fuse into one.
\end{walkbox}

\section{Looking back}

We started with dots and lines and ended at the algorithm that organizes the web, and every
step came from turning a network into a matrix and asking about its spectrum.

\begin{itemize}
\item The \textbf{adjacency matrix} $\mat{A}$ records connections; its powers count walks.
\item The \textbf{incidence matrix} $\mat{B}$ takes differences across edges -- a discrete
gradient -- and factors the Laplacian as $\mat{L} = \mat{B}\T\mat{B}$.
\item The \textbf{Laplacian} $\mat{L} = \mat{D}-\mat{A}$ is symmetric and positive
semidefinite; its energy is $\sum_{(i,j)\in E}(x_i-x_j)^2$; the multiplicity of its
eigenvalue $0$ counts connected components, and the next eigenvalue $\lambda_2$ (the Fiedler
value) measures connectivity and locates the best cut -- the foundation of spectral
clustering.
\item A \textbf{Markov chain} evolves distributions by $\vec\pi \mapsto \vec\pi\mat{P}$; its
long-run behavior is the left eigenvector of $\mat{P}$ for eigenvalue $1$, the stationary
distribution, guaranteed unique and reachable by Perron--Frobenius when the chain is
irreducible and aperiodic.
\item \textbf{PageRank} is the stationary distribution of a random surfer, made well-posed
by teleportation, and computed by the power method -- the same ``iterate and watch it
settle'' calculation, now recognized as eigenvector finding.
\end{itemize}

The unity behind the walking question is the lesson to carry away: connectivity, clustering,
mixing, and importance are all the same spectral story. Learn to read the eigenvalues of a
network's matrix and you can hear the shape of the network.

\subsection*{What We Can Now Do (Capabilities Unlocked)}
\begin{itemize}
  \item \textbf{Discrete Differential Geometry on Graphs:} Using the incidence matrix $\mat{B}$ and the graph Laplacian $\mat{L} = \mat{B}\T\mat{B}$, we generalize gradients, divergences, and Dirichlet energy ($\vec{x}\T\mat{L}\vec{x} = \sum (x_i-x_j)^2$) to discrete topological networks.
  \item \textbf{Spectral Partitioning and Cheeger Cuts:} The second smallest eigenvalue of the Laplacian ($\lambda_2$, the Fiedler value) and its corresponding eigenvector $\vec{v}_2$ enable spectral clustering, cutting graphs along bottleneck boundaries without checking an exponential number of combinatorial subsets.
  \item \textbf{Stationary Distribution Guarantees:} Using the Perron--Frobenius theorem, we analyze general stochastic transition matrices, proving when a Markov chain possesses a unique equilibrium distribution reachable from any arbitrary starting state.
  \item \textbf{Scalable Web-Scale Ranking:} By regularizing directed graphs with teleportation, we formulate global node importance as an eigenvector problem on a positive matrix ($\mat{G}\vec{\pi} = \vec{\pi}$), solved efficiently on billions of nodes via power iteration.
\end{itemize}

\subsection*{Honest Downsides and Limitations}
\begin{itemize}
  \item \textbf{Spectral Gap and Slow Mixing Bottlenecks:} The rate at which a Markov chain converges to its steady state is dictated by the spectral gap $1 - |\lambda_2|$. On graphs with severe bottlenecks (such as dumbbell graphs or barbell networks), $|\lambda_2| \approx 1$, causing mixing times to explode into millions of steps.
  \item \textbf{Stationarity Ignores Temporal Dynamics:} Markov chain steady states $\vec{\pi}$ assume static, time-invariant transition probabilities. They fail to capture real-world networks with bursty temporal edges, evolving network topologies, or memory-dependent higher-order interactions.
  \item \textbf{NP-Hardness of Optimal Graph Cuts:} While the Fiedler vector provides an effective heuristic for graph partitioning, Cheeger's inequality only guarantees an approximation; finding the true conductance-minimizing cut is NP-hard.
  \item \textbf{Memory Storage of Dense Teleportation:} Although the raw web graph is sparse, the regularized Google matrix $\mat{G} = \alpha\mat{P} + \frac{1-\alpha}{n}\vec{1}\vec{1}\T$ is completely dense ($n^2$ nonzeros). Forming $\mat{G}$ explicitly in RAM would be impossible; algorithms must exploit the low-rank rank-1 perturbation structure $\vec{1}\vec{1}\T$ to maintain $\mathcal{O}(|E| + n)$ iteration cost.
\end{itemize}

\subsection*{Where We Go Next in the Book}
\begin{itemize}
  \item \textbf{Chapter 36 (Optimization and Quadratic Forms):} Graph Dirichlet energy $\vec{x}\T\mat{L}\vec{x}$ is a symmetric positive semidefinite quadratic form. We generalize quadratic minimization to arbitrary symmetric systems, discovering how the Rayleigh quotient $\frac{\vec{x}\T\mat{A}\vec{x}}{\vec{x}\T\vec{x}}$ and the Conjugate Gradient algorithm find extrema without computing full eigenspaces.
  \item \textbf{Chapter 37 (Connections: Groups, Tensors, Functional):} We extend discrete graphs to continuous manifolds and continuous transformation groups, where graph Laplacians become the Laplace--Beltrami operator $\nabla^2$ of quantum mechanics and general relativity.
\end{itemize}

\section*{Exercises}
\addcontentsline{toc}{section}{Exercises}

\begin{exercise}[Build the three matrices]
\label{ch35:ex-build}
Consider the path graph on five nodes in a line: $1\!-\!2\!-\!3\!-\!4\!-\!5$.
\begin{enumerate}
\item Write down its adjacency matrix $\mat{A}$ and its degree matrix $\mat{D}$.
\item Write down the Laplacian $\mat{L} = \mat{D}-\mat{A}$ and verify that every row sums to
zero.
\item Orient each edge left-to-right and write the $4\times 5$ incidence matrix $\mat{B}$.
Confirm by direct multiplication that $\mat{B}\T\mat{B} = \mat{L}$.
\end{enumerate}
\end{exercise}

\begin{exercise}[Walks by hand]
\label{ch35:ex-walks}
For the running graph of Figure~\ref{ch35:fig-graph4}, we computed $\mat{A}^2$. Compute
$\mat{A}^3$ and verify that $(\mat{A}^3)_{11}$ equals the number of length-$3$ closed walks
from node $1$ back to node $1$. List those walks explicitly on the picture. (Hint: a
length-$3$ closed walk must use a triangle.)
\end{exercise}

\begin{exercise}[Counting components by eigenvalue]
\label{ch35:ex-components}
A graph on six nodes has two triangles, $\{1,2,3\}$ and $\{4,5,6\}$, with no edges between
them. Without computing all the eigenvalues, state the multiplicity of the eigenvalue $0$ of
its Laplacian and write down a basis for the corresponding eigenspace. Then add a single
edge $3\!-\!4$ joining the triangles; what is the new multiplicity of $0$, and why?
\end{exercise}

\begin{exercise}[The energy form]
\label{ch35:ex-energy}
Let $\mat{L}$ be the Laplacian of the running graph (Figure~\ref{ch35:fig-graph4}) and let
$\vec{x} = (1,1,-1,-1)\T$ (label nodes $1,2$ ``$+$'' and nodes $3,4$ ``$-$'').
\begin{enumerate}
\item Compute $\vec{x}\T\mat{L}\vec{x}$ two ways: directly as a matrix product, and as
$\sum_{(i,j)\in E}(x_i-x_j)^2$. They must agree.
\item Which edges contribute to the sum? Interpret the value as ``the number of edges cut''
by the partition $\{1,2\}$ versus $\{3,4\}$.
\end{enumerate}
\end{exercise}

\begin{exercise}[Stationary distribution]
\label{ch35:ex-stationary}
A frog hops between two lily pads, $A$ and $B$. From $A$ it stays with probability $0.7$ and
jumps to $B$ with probability $0.3$. From $B$ it stays with probability $0.4$ and jumps to
$A$ with probability $0.6$.
\begin{enumerate}
\item Write the $2\times2$ row-stochastic transition matrix $\mat{P}$.
\item Find the stationary distribution by solving $\vec\pi\mat{P} = \vec\pi$ with
$\pi_A + \pi_B = 1$.
\item Start at $\vec\pi^{(0)} = (1,0)$ and compute $\vec\pi^{(1)}, \vec\pi^{(2)},
\vec\pi^{(3)}$ to watch the convergence. What are the two eigenvalues of $\mat{P}$, and how
does the second one predict the speed you observe?
\end{enumerate}
\end{exercise}

\begin{exercise}[Why rows summing to one matters]
\label{ch35:ex-perron}
Prove that if $\mat{P}$ is row-stochastic (nonnegative entries, each row summing to $1$),
then $\mat{P}\vec{1} = \vec{1}$, so $1$ is an eigenvalue with right eigenvector $\vec{1}$.
Then explain why this does \emph{not} immediately give you the stationary distribution --
which eigenvector (left or right) do you actually need, and why are they different for a
non-symmetric $\mat{P}$?
\end{exercise}

\begin{exercise}[A reducible chain]
\label{ch35:ex-reducible}
Consider the chain on states $\{1,2,3\}$ with
\[
\mat{P} = \begin{bmatrix} 1 & 0 & 0\\ 0.5 & 0 & 0.5\\ 0 & 0 & 1 \end{bmatrix}.
\]
\begin{enumerate}
\item Draw the transition graph. Which states are ``absorbing'' (once entered, never left)?
\item Show that both $(1,0,0)$ and $(0,0,1)$ are stationary distributions. Why does this not
contradict Perron--Frobenius?
\item Starting from $(0,1,0)$, where does the chain end up, and how does the answer depend on
the coin flips? (This is why reducibility breaks uniqueness.)
\end{enumerate}
\end{exercise}

\begin{exercise}[PageRank by hand]
\label{ch35:ex-pagerank}
Use the four-page web of Figure~\ref{ch35:fig-web} with damping factor $\alpha = 0.85$.
\begin{enumerate}
\item Replace the dangling-free $\mat{P}$ given in the text and form the Google matrix
$\mat{G} = \alpha\mat{P} + \tfrac{1-\alpha}{4}\mat{J}$ explicitly (a $4\times4$ matrix).
\item Starting from uniform $\vec\pi^{(0)} = (\tfrac14,\tfrac14,\tfrac14,\tfrac14)$, run two
power-method steps $\vec\pi^{(0)}\mat{G}$ and $(\vec\pi^{(0)}\mat{G})\mat{G}$.
\item Rank the four pages by your $\vec\pi^{(2)}$. Does page $3$ come out on top? Does page
$4$ retain a small nonzero rank, and why is teleportation responsible?
\end{enumerate}
\end{exercise}

\begin{exercise}[The bridge between the two halves]
\label{ch35:ex-bridge}
For an undirected, connected graph let $\mat{L} = \mat{D}-\mat{A}$ be the Laplacian and
$\mat{P} = \mat{D}\inv\mat{A}$ be the random-walk transition matrix.
\begin{enumerate}
\item Verify the identity $\mat{P} = \mat{I} - \mat{D}\inv\mat{L}$.
\item Deduce that if $\vec{x}$ is an eigenvector of $\mat{D}\inv\mat{L}$ with eigenvalue
$\mu$, it is an eigenvector of $\mat{P}$ with eigenvalue $1-\mu$. In particular the
Laplacian's eigenvalue $0$ (constant eigenvector) maps to $\mat{P}$'s eigenvalue $1$. Explain
in words how this links ``connected $\Rightarrow$ unique constant ground state'' with
``connected $\Rightarrow$ unique stationary distribution.''
\end{enumerate}
\end{exercise}

\begin{exercise}[Fiedler in miniature, a proof]
\label{ch35:ex-fiedler}
Let $\mat{L}$ be the Laplacian of a graph with eigenvalues
$0 = \lambda_1 \le \lambda_2 \le \cdots \le \lambda_n$. Using the energy form
$\vec{x}\T\mat{L}\vec{x} = \sum_{(i,j)\in E}(x_i-x_j)^2$ and the fact (from the spectral
theorem for symmetric matrices) that
\[
\lambda_2 = \min_{\substack{\vec{x}\,\perp\,\vec{1}\\ \vec{x}\ne\vec0}}
\frac{\vec{x}\T\mat{L}\vec{x}}{\vec{x}\T\vec{x}},
\]
prove that $\lambda_2 = 0$ \emph{if and only if} the graph is disconnected. (One direction:
if disconnected, exhibit a nonconstant $\vec{x}\perp\vec1$ with zero energy. Other direction:
if $\lambda_2 = 0$, build a second independent constant-on-components vector and conclude
there are $\ge 2$ components.) This is Theorem~\ref{ch35:thm-components} in the
smallest nontrivial case.
\end{exercise}
